\documentclass[11pt]{article}
% Shared preamble for the hanzo-kernel Paper 07 series.
% One style, one vocabulary, inputted by every paper so the series reads as one voice.
\usepackage[margin=1in]{geometry}
\usepackage{booktabs}
\usepackage{amsmath}
\usepackage{amssymb}
\usepackage{graphicx}
\usepackage{xcolor}
\usepackage{hyperref}
\hypersetup{colorlinks=true, linkcolor=blue, urlcolor=blue, citecolor=blue}
\usepackage{enumitem}
\usepackage{listings}
\usepackage{array}

\definecolor{kw}{rgb}{0.13,0.29,0.53}
\definecolor{cmt}{rgb}{0.40,0.40,0.40}
\definecolor{str}{rgb}{0.16,0.50,0.16}
\lstset{
  basicstyle=\ttfamily\footnotesize,
  keywordstyle=\color{kw}\bfseries,
  commentstyle=\color{cmt}\itshape,
  stringstyle=\color{str},
  showstringspaces=false,
  breaklines=true,
  columns=fullflexible,
  frame=single,
  framesep=4pt,
}
\lstdefinelanguage{Rust}{
  morekeywords={fn,let,mut,pub,struct,impl,trait,enum,match,if,else,for,while,loop,return,use,mod,crate,self,super,async,await,where,type,const,static,unsafe,ref,move,dyn,as,true,false},
  morecomment=[l]{//}, morecomment=[s]{/*}{*/}, morestring=[b]", sensitive=true,
}

\newcommand{\x}{\ensuremath{\times}}
\newcommand{\dpa}{\texttt{dp4a}}
\newcommand{\hk}{\texttt{hanzo-kernel}}
\newcommand{\code}[1]{\texttt{#1}}

% Absorb minor prose overfulls without rivers; keep line-breaking sane.
\tolerance=800
\emergencystretch=3em


% --- Paper 03 local additions: theorem apparatus + notation ---
\usepackage{amsthm}
\newtheorem{theorem}{Theorem}
\newtheorem{proposition}{Proposition}
\newtheorem{lemma}{Lemma}
\newtheorem{corollary}{Corollary}
\theoremstyle{definition}
\newtheorem{definition}{Definition}
\theoremstyle{remark}
\newtheorem{remark}{Remark}

\newcommand{\R}{\mathbb{R}}
\newcommand{\E}{\mathbb{E}}
\newcommand{\mpi}{m_{\pi}}
\newcommand{\mJ}{m_{J}}
\newcommand{\Prob}{\mathcal{P}}
\newcommand{\Xspace}{\mathcal{X}}
\newcommand{\Jspace}{\mathcal{J}}

\title{\textbf{The Reward Is a Flock:\\ A Two-Population Mean-Field Adversarial Game for Fleet Training}}
\author{Hanzo AI --- ML Systems}
\date{hanzo-agentic Paper 03 --- the \emph{Continually Improving Fleets} series}

\begin{document}
\maketitle

\begin{abstract}
Reinforcement learning on language models is usually drawn as a policy optimizing against a reward. The
reward is treated as furniture --- a fixed function handed in from outside. This paper takes the reward as
the object of study and argues it should be neither fixed nor singular. A single learned judge is a
\emph{proxy}: optimize it hard enough and Goodhart's law guarantees the policy finds where the proxy and
the truth disagree. We replace the judge with a \emph{flock} --- a population of diverse judges whose
aggregate verdict is the reward --- and show three things. First, the flock is itself a mean-field system,
but run in the \emph{anti-ferromagnetic} regime: judges are rewarded for covering the judgment space, not
for agreeing, so the flock's stable configuration maximizes independence, and independence is exactly what
Condorcet's jury theorem needs and what shrinks the exploitable surface to the intersection of blind spots.
Second, the generator fleet (the ferromagnetic policy population of the companion paper) and the
discriminator flock couple through each other's mean fields with \emph{opposite sign}: the pair is a
two-population mean-field game, which is a generative adversarial network taken to the mean-field limit,
and its equilibrium is a Nash \emph{saddle} rather than a proxy peak. We make the anti-hacking claim
formal: a static judge field admits the reward-hacking maximizer as a stationary point, and a co-evolving
judge field does not, because the region the policy over-concentrates on is re-judged until the incentive
to stay there vanishes. Third, none of this is safe without an \emph{exogenous verifiable potential} ---
compile, test, benchmark, ship-and-stay-live --- a common term on both populations that breaks the
degenerate equilibrium in which flock and policy collude on a self-consistent but reality-detached reward.
We show the whole construction collapses, in the discrete per-group limit, onto the group-relative
advantage already computed by the shipped \texttt{hanzo-ml} GRPO trainer~\cite{shao2024grpo}: the flock's
anchored verdict is the per-group mean field, and GRPO's existing weighted multi-reward sum is the
non-adversarial, fixed-weight special case of a flock. We give an honest accounting of what is real in the
training substrate today and what the adversarial, anchored, co-evolving generalization still needs.
\end{abstract}

\section{The reward is the hard part}
\label{sec:intro}

A policy-gradient trainer is only as good as the number it climbs. The companion papers treat the
\emph{population} of policies: how a fleet aligns (Paper~01, \cite{hanzo-mfa}) and how it searches
(Paper~02, \cite{hanzo-evo}). Both assume a field to align to and a fitness to select on. This paper is
about that field. Where does it come from, and why is a single source of it dangerous?

The danger is Goodhart's law in its sharpest form. Let $q:\Xspace\to\R$ be the true quality of an output
$x$ (unknown, possibly undefinable pointwise) and let $h:\Xspace\to\R$ be a learned proxy for it. Optimizing
$\max_{x} h(x)$ does not return the $x$ that maximizes $q$; it returns the $x$ that maximizes $h-q$ plus
$q$, i.e.\ it seeks out and concentrates on the \emph{disagreement} between proxy and truth. A capable
policy is a very good optimizer, so it is a very good finder of that disagreement. This is reward hacking,
and it is not a bug in a particular reward model --- it is the generic behavior of maximizing a proxy.

Two structural moves defuse it, and this paper is about combining them. \textbf{Make the reward a
population} (a flock of judges, Section~\ref{sec:flock}), so that the exploitable disagreement must be
shared by many independent proxies at once. \textbf{Make the reward move} (co-evolve the flock against the
policy, Section~\ref{sec:twopop}), so that no static exploit is stationary. And underneath both,
\textbf{anchor the reward to something real} (a verifiable potential, Section~\ref{sec:anchor}), so the
flock cannot drift as a bloc.

\section{The flock, and why independence is the whole game}
\label{sec:flock}

\begin{definition}[Flock reward]
Let $\Jspace$ be a space of judges; a judge $j\in\Jspace$ maps an output to a scalar verdict
$j(x)\in\R$. Let $\mJ\in\Prob(\Jspace)$ be a probability measure over judges (the \emph{flock}). The flock
reward is a functional of the flock,
\[
  r(x;\mJ) \;=\; \mathrm{Agg}_{j\sim\mJ}\, j(x),
\]
with $\mathrm{Agg}$ the mean $\int j(x)\,d\mJ(j)$ in the simplest case, and a robust statistic (median,
trimmed mean, reliability-weighted mean) in practice.
\end{definition}

The mean case exposes the essential fact. Write judge $j$'s verdict as truth plus error,
$j(x)=q(x)+\varepsilon_j(x)$, with $\E[\varepsilon_j]=b$ (a shared bias) and
$\mathrm{Var}(\varepsilon_j)=\sigma^2$, and let $\bar\rho$ be the mean pairwise correlation of the
$\varepsilon_j$ across a flock of $K$ judges.

\begin{proposition}[Diversity is the reward's temperature]
\label{prop:condorcet}
The variance of the mean flock reward is
\[
  \mathrm{Var}\!\Big(\tfrac1K\textstyle\sum_{j=1}^K j(x)\Big)
  \;=\; \sigma^2\Big(\bar\rho + \tfrac{1-\bar\rho}{K}\Big)
  \;\xrightarrow[K\to\infty]{}\; \bar\rho\,\sigma^2 .
\]
Independent judges ($\bar\rho\to 0$) drive the variance to $0$; perfectly correlated judges
($\bar\rho\to 1$) leave it at $\sigma^2$ --- no gain from any number of them. The shared bias $b$ is
\emph{not} reduced by averaging at all; only decorrelated error is.
\end{proposition}

\begin{proof}
The variance of an average of $K$ unit-variance, mean-$\bar\rho$-correlated variables is
$\frac1{K^2}\big(K + K(K-1)\bar\rho\big)=\bar\rho+\frac{1-\bar\rho}{K}$; scale by $\sigma^2$. The bias term
is common to every judge and survives the average unchanged.
\end{proof}

Proposition~\ref{prop:condorcet} is the discriminative form of Condorcet's jury
theorem~\cite{condorcet1785}: a jury of weak but \emph{independent} voters converges to certainty, and a
jury of clones does nothing. Two consequences set the entire design.

\begin{corollary}[The exploitable surface is an intersection]
\label{cor:intersection}
A policy hacks a single judge $j$ by concentrating on $x$ where $\varepsilon_j(x)>0$. It hacks the flock
only on $x$ where $\varepsilon_j(x)>0$ for enough $j$ simultaneously --- a set that shrinks with judge
diversity, and vanishes for the components of error that are independent across the flock. Only the shared
bias $b$ (the common blind spot) remains hackable; Section~\ref{sec:anchor} is how we remove it.
\end{corollary}

\begin{remark}[Judge diversity is Paper~01's temperature]
Independence is not free: it must be produced. Judges drawn from one model family, one rubric, one
persona, one decoding temperature are near-clones ($\bar\rho\to1$). Diversity across model families,
prompts, rubrics, and sampling temperatures is what keeps $\bar\rho$ low --- and $\bar\rho$ plays exactly
the role of the alignment order parameter $m$ in Paper~01~\cite{hanzo-mfa}. A flock that has aligned into
consensus is a flock frozen below its Curie point: coherent, and useless.
\end{remark}

\section{The flock is an anti-ferromagnet}
\label{sec:antiferro}

If Paper~01 asks the policy population to align --- a ferromagnetic mean field with coupling $J>0$, whose
stable state is consensus --- then the flock must want the opposite. We give the judges a mean-field
interaction with \emph{negative} coupling.

\begin{definition}[Anti-ferromagnetic flock objective]
Each judge minimizes a cost with three terms,
\[
  c_J(j;\mJ,V)
   \;=\; \underbrace{\alpha\,\ell_{\text{ver}}\!\big(j,V\big)}_{\text{agree on the verifiable}}
   \;-\; \underbrace{|J|\!\!\int\! K(j,j')\,d\mJ(j')}_{\text{repel on the subjective, }J<0}
   \;+\; \underbrace{\lambda\,\Omega(j)}_{\text{regularizer}} ,
\]
where $\ell_{\text{ver}}$ penalizes disagreement with the verifiable potential $V$ of
Section~\ref{sec:anchor}, $K(j,j')$ is a similarity kernel on judges, and the middle term \emph{rewards}
being far from the rest of the flock (a repulsive, $J<0$, mean-field interaction).
\end{definition}

The flock's equilibrium is therefore \emph{structured}: ferromagnetic (agreeing, $J>0$ effectively) on
the axis $V$ can measure, and anti-ferromagnetic (diverse, $J<0$) on the subjective axis it cannot. This is
the correct anisotropy. On ``does it compile,'' we want the judges to agree, because there is a truth to
agree on. On ``is the abstraction clean,'' we want them to spread, because the aggregate of diverse
subjective reads is a better estimator than any one, and because a spread flock is unhackable where a
consensus flock is a single point of failure. The repulsive term is what actively \emph{produces} the low
$\bar\rho$ that Proposition~\ref{prop:condorcet} needs; diversity stops being a hope and becomes an
optimized quantity.

\section{Two populations, opposite signs: the mean-field GAN}
\label{sec:twopop}

Now couple them. Let $\mpi\in\Prob(\Xspace)$ be the policy population's output measure (the generator) and
$\mJ\in\Prob(\Jspace)$ the flock (the discriminator). Each population runs a mean-field game --- a
Hamilton--Jacobi--Bellman equation for the individual's optimal behavior against the population, coupled to
a Fokker--Planck--Kolmogorov equation for how the population evolves under that behavior~\cite{lasry2007mfg}
--- and the two are coupled through each other's mean fields.

The generator's individual cost \emph{falls} with reward:
$c_\pi(x;\mJ) = -\,r(x;\mJ) + \tfrac1\beta\,\mathrm{KL}(\pi\,\|\,\pi_{\text{ref}})$, where the KL-to-reference
term is the temperature $1/\beta$ (Paper~01). The flock's cost \emph{rises} with the reward's error against
truth: it is paid to make $r(\cdot;\mJ)$ track $q$ (anchored by $V$) and to deny the generator any free
lunch. The two objectives meet with opposite sign in $r$: the generator pushes $\mpi$ toward high $r$; the
flock moves $\mJ$ to keep $r$ honest exactly where $\mpi$ has gone.

\begin{definition}[Two-population adversarial MFG]
An equilibrium is a pair $(\mpi^\ast,\mJ^\ast)$ such that $\mpi^\ast$ is a mean-field best response to
$\mJ^\ast$ in the generator game and $\mJ^\ast$ is a mean-field best response to $\mpi^\ast$ in the flock
game. Equivalently it is a saddle of the min--max functional
\[
  \min_{\mJ}\ \max_{\mpi}\ \; \E_{x\sim\mpi}\big[\,r(x;\mJ)\,\big]\;-\;\gamma\, D\big(r(\cdot;\mJ),\,q\big),
\]
the generator maximizing the reward it is scored by, the flock minimizing the divergence $D$ of that reward
from truth.
\end{definition}

This is a generative adversarial network~\cite{goodfellow2014gan} written at the mean-field limit:
generator population vs.\ discriminator population, each a distribution rather than a single network. Its
value is a saddle, not a maximum --- and that distinction is the entire anti-hacking argument.

\begin{theorem}[Existence under monotone coupling]
\label{thm:existence}
Suppose the running costs are convex and Lipschitz in the control, the coupling functionals
$\mpi\mapsto c_J$ and $\mJ\mapsto c_\pi$ are continuous for the Wasserstein topology and satisfy the
Lasry--Lions monotonicity condition, and the state spaces are compact. Then the two-population system
admits at least one equilibrium $(\mpi^\ast,\mJ^\ast)$.
\end{theorem}

\begin{proof}[Proof sketch]
Each single-population game is a standard MFG whose best-response map is well defined and continuous under
the stated assumptions~\cite{lasry2007mfg}. The joint best-response
$(\mpi,\mJ)\mapsto(\mathrm{BR}_\pi(\mJ),\mathrm{BR}_J(\mpi))$ is then a continuous self-map of the product
of two convex, compact Wasserstein-ball measure sets; Schauder's fixed-point theorem gives a fixed point,
which is by definition an equilibrium. Monotonicity gives uniqueness in the standard one-population
regularity class; we do \emph{not} claim uniqueness here, because the adversarial (opposite-sign) coupling
is generically \emph{non}-monotone --- see Remark~\ref{rem:limitcycle}.
\end{proof}

\begin{remark}[Honest boundary: non-monotone coupling can cycle]
\label{rem:limitcycle}
The opposite-sign coupling that makes this a game rather than a potential-minimization is exactly what
breaks Lasry--Lions monotonicity. Without the anchor of Section~\ref{sec:anchor}, the flow
$(m_{\pi,t},m_{J,t})$ can limit-cycle --- the generator chases the reward, the flock chases the generator,
neither settles --- the mean-field echo of GAN training instability. The verifiable potential is what turns
the cycle into a pinned saddle; it is not decoration, it is the well-posedness.
\end{remark}

\section{Why a moving reward cannot be hacked}
\label{sec:saddle}

Here is the payoff, stated as sharply as it deserves.

\begin{proposition}[Static field admits the hack; co-evolving field does not]
\label{prop:hack}
Fix the flock at $\mJ^0$, so the reward is a static field $h(x)=r(x;\mJ^0)$. The generator's problem
$\max_{\mpi}\E_{\mpi}[h]$ is a linear program on $\Prob(\Xspace)$; its optimum places all mass on
$\arg\max_x h(x)$. If $x^\star\in\arg\max h$ has $q(x^\star)<h(x^\star)$ --- a proxy--truth gap, which by
Corollary~\ref{cor:intersection} exists wherever the flock's shared bias is nonzero --- then the reward-hacking
concentration on $x^\star$ is the \emph{stationary point} of the generator dynamics.

Now let the flock co-evolve: $\mJ=\mJ[\mpi]$ is the flock's best response to the current policy
population. The generator faces $\max_{\mpi}\E_{\mpi}\big[r(x;\mJ[\mpi])\big]$, in which the field is a
functional of the very measure being optimized. The former stationary point $x^\star$ is \emph{not}
stationary: as $\mpi$ concentrates on $x^\star$, the flock's best response (i) increases judge coverage
near $x^\star$ (the anti-ferromagnetic repulsion drives judges into the newly dense, newly important
region) and (ii) is pulled by the verifiable term toward $V(x^\star)$, which is low. Both lower
$r(x^\star;\mJ[\mpi])$ until the generator's incentive to remain at $x^\star$ vanishes. No pure exploit is
a fixed point of the coupled system; the only fixed points are saddles where the generator can no longer
improve the reward without the flock immediately correcting it.
\end{proposition}

The one-line reading: \emph{a target that re-aims at wherever you aimed last cannot be shot with a fixed
aim.} Reward hacking is optimization against a stationary proxy; remove the stationarity and you remove the
hack. This is the formal content of ``rotate the judges,'' ``adversarial judges,'' and ``the discriminator
must keep up with the generator,'' unified as one property of the coupled equilibrium.

\section{The anchor: a verifiable potential on both populations}
\label{sec:anchor}

A flock of \emph{AI} judges has a failure mode that no amount of diversity fixes: it can drift as a bloc.
If the shared bias $b$ of Proposition~\ref{prop:condorcet} is itself learned and shared --- a common
sycophancy, a common taste, a hallucinated notion of quality --- then the flock and the policy can settle
into a self-consistent equilibrium that is internally coherent and externally false. In game terms, the
min--max of Section~\ref{sec:twopop} without an external term has a degenerate family of equilibria: any
mutually-consistent $(\mpi,\mJ)$ pair, including collapsed and reality-detached ones.

\begin{definition}[Verifiable potential]
Let $V:\Xspace\to\R$ be an \emph{exogenous}, non-learned reward computable by execution: the code compiles;
the tests pass; the benchmark clears its bit-exact gate; the change ships and stays live. $V$ is a common
potential added to \emph{both} populations --- it is a term in the flock's $\ell_{\text{ver}}$ (judges are
penalized for verdicts inconsistent with $V$) and a term in the generator's reward.
\end{definition}

\begin{proposition}[Anchoring breaks the degeneracy]
\label{prop:anchor}
With $V$ present, every equilibrium $(\mpi^\ast,\mJ^\ast)$ must be consistent with $V$ on the verifiable
axis: the flock's $V$-inconsistent verdicts carry strictly positive cost, so no equilibrium can place mass
on outputs that fail $V$ while claiming high reward. The set of equilibria is pinned to the
$V$-consistent manifold, eliminating the collapsed/hallucinated family and turning
Remark~\ref{rem:limitcycle}'s cycle into a bounded saddle.
\end{proposition}

This is the RLVR (reinforcement learning from verifiable rewards) anchor married to RLAIF (reinforcement
learning from AI feedback~\cite{bai2022constitutional}): the \emph{verifiable} axis is judged by execution, ferromagnetically (there is
a truth, the judges should agree with it and each other); the \emph{subjective} axis is judged by the
anti-ferromagnetic flock. $V$ does not have to be dense or complete --- it only has to be \emph{real} and
\emph{present}, enough to make reality a term the equilibrium cannot ignore.

\begin{remark}[Judges graduate through the gate]
The co-evolution is closed by promoting judges: a fine-tuned enso model may \emph{join} the flock, but only
after clearing the verifiable gate on a held-out battery, and it is admitted with a reliability weight
proportional to its historical agreement with $V$. A judge that keeps losing to the anchor is down-weighted.
This is how the flock improves without eating its own tail: new judges enter through reality, not through
consensus.
\end{remark}

\section{Both populations live at criticality; the GA moves them}
\label{sec:crit}

Paper~01 argues each population should sit just below its ordering transition --- coherent enough to
converge, hot enough to correct. That applies to both here: a generator fleet frozen into one output has no
exploration, and a flock frozen into one opinion has no robustness. The control knobs are the same
temperatures --- decoding temperature and KL coefficient for the generator, judge diversity ($1-\bar\rho$)
and coverage weight for the flock --- annealed over training, never quenched. Paper~02's genetic
search~\cite{hanzo-evo} operates on \emph{both} populations: it evolves the policy variants (crossover of
partial solutions) and it evolves the flock (recombine rubrics, retire drifted judges, breed diverse new
ones). The three papers are therefore one coupled master system: a \textbf{generator fleet}
$\otimes$ \textbf{discriminator flock}, \textbf{searched} by a genetic algorithm, \textbf{anchored} by a
verifiable potential, all closing on a single mean-field number.

\section{The grounding: it collapses onto GRPO}
\label{sec:grounding}

The construction above is not a proposal to build a new trainer. It is the adversarial, anchored,
co-evolving generalization of the loop the \texttt{hanzo-ml} stack already runs. The shipped GRPO
trainer~\cite{hanzo-native-training,shao2024grpo} samples a \emph{group} of $G$ rollouts per prompt, scores
each, and forms the group-relative advantage
\[
  A_i \;=\; \frac{r_i - \mathrm{mean}_{k\le G}\, r_k}{\mathrm{std}_{k\le G}\, r_k}
  \qquad(\texttt{math::batch\_advantages},\ \text{ddof}=1),
\]
then a length-normalized policy-gradient loss with an optional $k3$ KL penalty to a frozen reference. Read
this against Sections~\ref{sec:flock}--\ref{sec:twopop}:

\begin{itemize}[leftmargin=1.4em,itemsep=2pt]
  \item \textbf{$r_i$ is the anchored flock verdict} on rollout $i$: the aggregate of the diverse judges
        plus the verifiable term $V$. GRPO's existing ``weighted multi-reward sum'' (\texttt{score()},
        \texttt{trainer.rs:119}) is precisely a flock with \emph{fixed weights and no adversary} --- the
        non-adversarial special case. Turning those fixed weights into a co-evolving, diversity-optimizing,
        $V$-anchored flock is the whole delta.
  \item \textbf{The group mean is the mean field.} $\mathrm{mean}_{k}\,r_k$ is the finite-sample estimator
        of $\E_{\mpi}[r]$; subtracting it is the mean-field coupling of Paper~01 computed per group. The
        $\mathrm{std}$ normalization is the local temperature. GRPO is, structurally, a per-group
        discretization of the mean-field game.
  \item \textbf{The $k3$ KL is $1/\beta$.} The frozen-reference penalty is the generator temperature that
        Section~\ref{sec:twopop} keeps at criticality.
\end{itemize}

\paragraph{Cost.} A flock of $K$ judges is $K\times$ the eval cost per rollout. Two mitigations, both
structural. \emph{Cascade}: score with one cheap judge first; escalate to the full flock only on high stakes
or when the cheap judge's uncertainty (or disagreement with $V$) is high. \emph{Uncertainty gating}: the
flock spread is an epistemic-uncertainty estimate (Proposition~\ref{prop:condorcet}); train hard where the
flock agrees and $V$ confirms, and treat high-spread rollouts as exploration signal rather than gradient.

\paragraph{Where it lives.} The eval-to-finetune bridge is the \texttt{/v1/evals} product (Hanzo's
LLM-as-a-judge / observability plane): the flock scores every output there, the verifiers anchor it, and
the aggregate is emitted as the reward stream the GRPO loop consumes. The judge population, its diversity
objective, and its recalibration schedule are configuration on that plane, not new infrastructure.

\section{What is real, what needs building, what is out of scope}
\label{sec:accounting}

\paragraph{Real (in the substrate today).} The GRPO five-phase loop --- rollout, weighted multi-reward,
group-relative advantage with $\text{ddof}=1$, length-normalized policy-gradient loss, $k3$ KL, AdamW step ---
is ported and tested~\cite{hanzo-native-training}. The verifiable signals ($V$) exist as CI / harness
outcomes (compile, test, benchmark gates, ship-and-live telemetry). The \texttt{/v1/evals} plane hosts
LLM-as-a-judge scoring. Multi-model consensus primitives exist in the tool layer.

\paragraph{Needs building.} The judge population as a first-class object: a diversity (anti-ferromagnetic)
objective that actively minimizes $\bar\rho$; the co-evolution loop (flock best-responds to the policy
population, drifted judges down-weighted, graduated judges admitted through the $V$ gate); the cascade /
uncertainty-gated escalation; and the wiring that turns the flock's anchored verdict into the per-group
reward stream on \texttt{/v1/evals}. Today's ``weighted multi-reward'' is the fixed-weight, non-adversarial
stub of the flock.

\paragraph{Out of scope.} A full Lasry--Lions well-posedness proof for this specific non-monotone coupling
(we invoke existence under monotonicity, Theorem~\ref{thm:existence}, and are explicit that the adversarial
coupling is non-monotone and can cycle without the anchor, Remark~\ref{rem:limitcycle}); learned
\emph{optimal} judge weights and kernels; and any human-in-the-loop calibration schedule beyond the
verifiable gate.

\section{Conclusion}
\label{sec:conclusion}

The reward is not furniture. A single judge is a stationary proxy and will be hacked; a flock of diverse
judges shrinks the exploitable surface to the intersection of blind spots
(Proposition~\ref{prop:condorcet}, Corollary~\ref{cor:intersection}), a flock that repels into diversity
(anti-ferromagnetic, Section~\ref{sec:antiferro}) actively produces the independence that shrinkage needs,
a flock that co-evolves with the policy removes the stationarity that hacking requires
(Proposition~\ref{prop:hack}), and a verifiable potential pins the whole adversarial equilibrium to reality
so the flock cannot drift as a bloc (Proposition~\ref{prop:anchor}). Structurally it is a two-population
mean-field game --- a GAN at the mean-field limit --- and it collapses, per rollout group, onto the
group-relative advantage the \texttt{hanzo-ml} GRPO trainer already computes, with today's fixed-weight
multi-reward as its non-adversarial special case. Together with the fleet's ferromagnetic alignment
(Paper~01) and the genetic search over both populations (Paper~02), it completes one system:
\emph{compass, engine, thermostat} --- generator fleet $\otimes$ discriminator flock, searched, anchored,
and closing on a single mean field.

\begin{thebibliography}{10}

\bibitem{hanzo-mfa}
Hanzo AI --- ML Systems. (2026).
\textit{Critical Alignment of Agentic Fleets: Mean-Field Games and Ferromagnetic Order}.
hanzo-agentic Paper~01, the \emph{Continually Improving Fleets} series.

\bibitem{hanzo-evo}
Hanzo AI --- ML Systems. (2026).
\textit{Evolutionary Orchestration: Population Search over Agentic Fleets}.
hanzo-agentic Paper~02, the \emph{Continually Improving Fleets} series.

\bibitem{hanzo-native-training}
Hanzo AI --- ML Systems. (2026).
\textit{Native Training in hanzo-ml/hanzo-engine: One Engine, One Quant Format, One Backend}.
Hanzo Industries Research.

\bibitem{hanzo-consensus-zap}
Hanzo AI --- ML Systems. (2026).
\textit{The Three-Plane Compute Fabric: Consensus and ZAP Compute}.
\emph{One Source, Every Backend} series.

\bibitem{shao2024grpo}
Shao, Z., Wang, P., Zhu, Q., Xu, R., Song, J., et al. (2024).
\textit{DeepSeekMath: Pushing the Limits of Mathematical Reasoning in Open Language Models}.
arXiv:2402.03300. (Group Relative Policy Optimization.)

\bibitem{lasry2007mfg}
Lasry, J.-M., \& Lions, P.-L. (2007).
\textit{Mean Field Games}.
Japanese Journal of Mathematics, 2(1), 229--260.

\bibitem{condorcet1785}
de Condorcet, N. (1785).
\textit{Essai sur l'application de l'analyse \`a la probabilit\'e des d\'ecisions rendues \`a la pluralit\'e des voix}.
Imprimerie Royale, Paris. (The jury theorem.)

\bibitem{goodfellow2014gan}
Goodfellow, I., Pouget-Abadie, J., Mirza, M., Xu, B., Warde-Farley, D., Ozair, S., Courville, A., \& Bengio, Y. (2014).
\textit{Generative Adversarial Networks}.
NeurIPS 2014.

\bibitem{skalse2022reward}
Skalse, J., Howe, N.~H.~R., Krasheninnikov, D., \& Krueger, D. (2022).
\textit{Defining and Characterizing Reward Hacking}.
NeurIPS 2022.

\bibitem{bai2022constitutional}
Bai, Y., Kadavath, S., Kundu, S., et al. (2022).
\textit{Constitutional AI: Harmlessness from AI Feedback}.
arXiv:2212.08073.

\end{thebibliography}

\end{document}
